% ==================================================================================== %
\documentclass{article}
\usepackage[a4paper]{geometry}
\usepackage[english]{babel}
\usepackage{parskip}
\usepackage{amsmath, amssymb}
\usepackage{graphicx}
\usepackage{subcaption}
\usepackage{float}
\usepackage[hidelinks]{hyperref}
\usepackage{listings}
\usepackage{xcolor}
% ==================================================================================== %
\definecolor{nord-bg}{HTML}{FFFFFF}
\definecolor{nord-fg}{HTML}{2E3440}
\definecolor{nord-comment}{HTML}{5E81AC}
\definecolor{nord-keyword}{HTML}{BF616A}
\definecolor{nord-string}{HTML}{A3BE8C}
\definecolor{nord-number}{HTML}{B48EAD}
\definecolor{nord-operator}{HTML}{D08770}
% ==================================================================================== %
\lstdefinestyle{nord_light}{
  backgroundcolor=\color{nord-bg},
  basicstyle=\color{nord-fg}\footnotesize\ttfamily,
  commentstyle=\color{nord-comment},
  keywordstyle=\color{nord-keyword}\bfseries,
  stringstyle=\color{nord-string},
  numberstyle=\color{nord-number},
  otherkeywords={self,cls,@classmethod,@property},
  keywordstyle={[2]\color{nord-operator}},
  frame=leftline,
  showstringspaces=false,
  breaklines=true,
  tabsize=4,
  extendedchars=true,
  numbers=left,
  numberstyle=\tiny\color{nord-fg},
  stepnumber=1,
  numbersep=10pt,
}
\lstset{style=nord_light}
% ==================================================================================== %
\title{NUR A - Assignment 2}
\author{Taotao Yang (s4866835)}
\date{Dated: March 14, 2026} %\today
% ==================================================================================== %
\begin{document}
\maketitle
% ==================================================================================== %

Write a short description of your solution in each section. 
Explain your choices. Explicitly draw conclusions from your results.

You can refer to selected lines of code where you think it is necessary. 
However, you must provide the full code relevant to each (sub)question below it.

Don't forget to describe your figures and mention what conclusions you draw from them!

% ==================================================================================== %
\section{Satellite galaxies around a massive central}

  \subsection{(1a) Determining the normalization}

    Explain the integration method used (e.g. Romberg integration or another Richardson 
    method). Describe how you implemented the spherical volume element \(dV\).

    \begin{table}[H]
    \centering
    \caption{Normalization constant of the satellite galaxy profile.}
      \begin{tabular}{|c|}
        \hline
        Normalization constant \(A\) \\
        \hline
        \input{output/a2q1_satellite_A.txt} \\
        \hline
      \end{tabular}
    \end{table}

  \subsection{Code for the integrator}
    \lstinputlisting[language=Python]{output/a2q1_satellites_integrator_code.txt}

  \subsection{(1b) Sampling the satellite distribution}
    Describe the sampling method used (e.g. rejection sampling). Explain how the random 
    number generator was implemented and how its randomness was tested.

    \begin{figure}[H]
      \centering
      \includegraphics[width=0.85\textwidth]{plots/a2q1_my_solution_1b.png}
      \caption{
        Comparison between the analytical satellite distribution \(N(x)\) and 
        the histogram of 10 000 sampled satellite galaxies in logarithmic bins.
      }
    \end{figure}

    Does your sampled distribution agree with the analytical prediction?

  \subsection{Code for sampling}
    \lstinputlisting[language=Python]{output/a2q1_satellites_sampling_code.txt}
 
  \subsection{(1c) Selecting and sorting satellites}
    Describe the selection of 100 galaxies (equal probability, no duplicates, no 
    rejection). Explain the sorting algorithm you used.

    \begin{figure}[H]
      \centering
      \includegraphics[width=0.8\textwidth]{plots/a2q1_my_solution_1c.png}
      \caption{Cumulative number of selected galaxies as a function of relative radius.}
    \end{figure}

  \subsection{Code for selection and sorting}

    \lstinputlisting[language=Python]{output/a2q1_satellites_selection_code.txt}

  \subsection{(1d) Numerical derivative}

    \begin{table}[H]
    \centering
    \caption{Derivative of the satellite number density at \(x=1\).}
      \begin{tabular}{|c|c|}
        \hline
        Method & Value \\
        \hline
        Analytical & \input{output/a2q1_satellite_deriv_analytic.txt} \\
        Numerical & \input{output/a2q1_satellite_deriv_numeric.txt} \\
        \hline
      \end{tabular}
    \end{table}

    Do your analytical and numerical derivatives agree? Why?

  \subsection{Code for derivative calculation}
    \lstinputlisting[language=Python]{output/a2q1_satellites_derivative_code.txt}

\section{Heating and cooling in HII regions}

  \subsection{(2a) Photoionization heating and radiative recombination cooling}
    Describe the implementation of your root-finding method (e.g. bisection, 
    Newton-Raphson).

    \begin{table}[H]
    \centering
    \caption{Equilibrium temperature for the simplified heating-cooling model.}
      \begin{tabular}{|c|c|c|}
        \hline
        Temperature [K] & Absolute error & Relative error \\
        \hline
        \input{output/a2q2_equilibrium_temp_simple.txt}
        \\
        \hline
      \end{tabular}
    \end{table}

    How many iterations it took? Discuss the convergence of your root finder.

  \subsection{(2b) More realistic model}

    \begin{table}[H]
      \centering
      \caption{Equilibrium temperature for different gas densities.}
      \begin{tabular}{|c|c|}
        \hline
        Density \(n_H\) [cm\(^{-3}\)] & Equilibrium temperature [K] \\
        \hline
        $10^{-4}$ & \input{output/a2q2_equilibrium_low_density.txt} \\
        $1$ & \input{output/a2q2_equilibrium_mid_density.txt} \\
        $10^4$ & \input{output/a2q2_equilibrium_high_density.txt} \\
        \hline
      \end{tabular}
    \end{table}

    Discuss the results. What is the difference for the low, intermediate and high gas 
    density?

  \subsection{Code for Exercise 2}
    \lstinputlisting[language=Python]{output/a2q2_heatingcooling_code.txt} 

\section{Conclusions}

\end{document}
